\chapter*{前言}

这本书从硬件本身开始。一根线上连续变化的电压怎样被稳定解释成 0 和 1；一个受控开关怎样让一个信号决定另一条电流路径是否导通；多个开关怎样组成逻辑门、寄存器、ALU 和控制器——只有这些底层约定成立，CPU 和整机接口才具有可靠含义。

全书分为七个部分、二十六章。第一部分从电平、器件、组合逻辑和时序逻辑走到 RTL；第二部分建立数值、ISA、数据通路、流水线、异常和性能的处理器主线；第三部分用七章完整展开存储层次、Cache、地址转换、DRAM、SSD、硬盘和存储 I/O；第四部分讲总线、PCIe 与设备事务；第五部分讲电源、时钟、主板、固件和平台控制；第六部分讲外设、网络、GPU 以及链接装载；第七部分把所有边界收束到一台可搭建、可观察、可排错的整机。

CPU 在本书中是前述部件按同步边界组织后的结果，不是作为孤立名词提前给出的。存储也不被当作一个抽象方框：每一层都要说明物理单元、阵列、控制器、接口协议、访问粒度、完成语义和错误路径。读完整本书后，读者应能把一台机器从电平、门电路、时序逻辑、运算部件、数据通路和控制器，一路连到主存、持久化存储、互连、设备、固件与系统软件边界。

书中的训练以纸面推演、波形阅读、规格分析和 trace 走读为主。有 FPGA 或开发板的读者可以把任务落实到真实硬件，但没有硬件也能完成全部练习。最终目标不是记住部件名称，而是面对一次取指、访存、I/O 或故障时，能够指出数据从哪里来、由谁控制、何时提交、通过什么接口完成，以及失败后从哪里开始定位。
